\chapter{10 m wind and gust}
\label{cap:raffica}

\section{The \texorpdfstring{\campo{10fg}}{10fg} field}
\label{sec:10fg}

The \campo{10fg} field (paramId 49) is the maximum \SI{10}{m} wind speed that the simulation
produced in the hour preceding the validity time. It comes from \campo{WSPD10MAX}, one of the hourly
maxima that WRF writes with \campo{nwp_diagnostics=1}. It is a \emph{resolved wind}: it contains no
turbulent term and no estimate of what the grid does not represent \parencite{issue50}.

The ERA5 gust is a different quantity. The World Meteorological Organization defines the gust as
the maximum of the wind averaged over \SI{3}{s}, an interval shorter than one time step of a
\SI{31}{km} model, and ERA5 publishes an estimate of it, the \emph{parameterized gust}
\parencite{paramdb}. In the IFS the parameterized gust is the mean \SI{10}{m} wind plus a turbulent
term proportional to the friction velocity, plus, in convective conditions, a term proportional to
the wind shear in the \SIrange{850}{950}{\hPa} layer \parencite[sec.~3.10.4, eq.~3.99]{ifs41r2p4}. The
IFS computes it at every time step and archives its maximum since the last output.

The regional simulation receives from ERA5 the state of the atmosphere and of the surface, not the
diagnostic fields. The standard ECMWF field table of WPS 4.1 (\campo{Vtable.ECMWF}) contains
neither paramId 49 nor 228029 \parencite{wps41}; the table used by this simulation has not been
verified. Among the 200 variables of a wrfout file no gust field is present besides
\campo{WSPD10MAX}, and the field published as \campo{10fg} is the closest to a gust that the
simulation provides.

\section{Comparison with ERA5}
\label{sec:10fg-era5}

\Cref{tab:10fg} gives the comparison over 16 validity times. \campo{10fg} is \num{0.625} times the
ERA5 \campo{10fg} and \num{0.644} times the instantaneous parameterized gust \campo{i10fg}, with
correlations of \num{0.813} and \num{0.820}. Against the instantaneous \SI{10}{m} wind of the same
file it is between \num{1.004} and \num{1.032} times that wind, and at none of the \num{2220273} points is it
lower. Under paramId 49 the archive publishes the \SI{10}{m} wind speed, whereas the ERA5
parameterized gust is about \num{1.6} times its mean wind. The ratio to the ERA5 \campo{10fg} stays
between \num{0.593} and \num{0.673} at all 34 validity times of the verification sample, in every
month and at every hour, and the difference is one of definition.

\begin{table}[htbp]
  \centering
  \small
  \caption{The archive's \campo{10fg} against three references, over 16 validity times.}
  \label{tab:10fg}
  \begin{tabular}{@{}lll@{}}
    \toprule
    Reference & CHAPTER / reference & Correlation \\
    \midrule
    ERA5 \campo{10fg} (parameterized gust, hourly maximum) & \num{0.625} & \num{0.813} \\
    ERA5 \campo{i10fg} (parameterized gust, instantaneous) & \num{0.644} & \num{0.820} \\
    CHAPTER instantaneous \SI{10}{m} wind, same file & \numrange{1.004}{1.032} & \\
    \bottomrule
  \end{tabular}
\end{table}

At \SI{3}{km} part of what ERA5 has to parameterize is resolved, as the friction velocity over land
shows, between \num{1.62} and \num{2.06} times that of ERA5 (\cref{sec:rugosita}). Read as a gust,
\campo{10fg} underestimates the gust by about \SI{38}{\percent}, and no bias correction turns it into a gust,
because it is a different quantity. With a correlation of \num{0.813} with the ERA5 \campo{10fg},
the field is suited to ranking places and hours by windiness and to describing the spatial
structure of the wind, and it acts as a lower bound on the gust; it is not suited to estimating the
exceedance of a gust threshold.

\section{The hourly window}
\label{sec:10fg-finestra}

The \campo{10fg} messages declare a maximum over the preceding hour (\cref{sec:tempo}). The hourly
ERA5 parameterized gust is also a maximum over the preceding hour, so the length of the window does
not explain the difference between the two fields. The encoding is correct. WRF resets \campo{WSPD10MAX} each
time it writes a wrfout file, but over six whole days, 15 July 2019 and 1 and 15 February,
20 March, 15 July and 15 August 2024, for 138 transitions from one hour to the next, the reset was
missed 17 times, \SI{12.3}{\percent} of the cases \parencite{issue49}. In those hours the field
contains the maximum over two hours or more, under a message that declares one.

The effect shows in the mean excess over the instantaneous \SI{10}{m} wind:
\SIrange{0.068}{0.311}{m/s} in the hours with a reset and \SIrange{0.499}{0.784}{m/s} in those
without. On a domain mean of about \SI{4.5}{m/s} this is an overestimate of
\SIrange{10}{17}{\percent}. In the hours without a reset the field approaches the ERA5
parameterized gust because the window is longer, and the value alone cannot identify them. Nor do the other fields of the file: the seven WRF hourly maxima are reset by the same
event and, when the reset is missed, all carry over together, so none of them signals the missed
reset of the others. Recognizing those hours requires the wrfout file of the preceding hour, which
is not available at conversion time. The mechanism of the missed reset has not been established.

In statistics over a season, \SI{12.3}{\percent} of the hours carry a window longer than the
declared one; for a single hour the declared window of \campo{10fg} is not guaranteed.
